\section{Alucinaciones y jailbreak: cómo compiten varios circuitos}

El cómputo en paralelo no solo aporta capacidades, también genera conflictos. El modelo puede identificar una entidad, preparar una respuesta y activar una tendencia de «no sé» o de refusal al mismo tiempo; la salida final depende de qué rutas se inhiben y cuáles son más fuertes. Esta sección explica la hallucination y el jailbreak como mecanismos de competencia local, no como un interruptor global.

Al pasar de la suma a la conducta de seguridad, primero hay que cambiar el criterio de lectura de las figuras: en las figuras de aritmética el token objetivo tiene un valor verdadero definido, mientras que «si se debe responder, rechazar o expresar incertidumbre» depende además de la cobertura de conocimiento, de las políticas y del contexto. Por eso esta sección compara primero pares de prompts, luego aplica feature suppression y, por último, solo plantea hipótesis de mecanismo de alcance limitado. Nos interesa qué evidencia interna cambió la tendencia de la salida; no equiparamos directamente el nombre de una feature con la truthfulness o la safety del sistema completo.

\lecturefigure{slide-44-hallucinations-intro.jpg}{Parallel Processing: Hallucinations}{00:59:56}

\subsection{Una tendencia IDK always-on}

Comparemos a Michael Jordan con el ficticio Michael Batkin: el primero activa evidencia de known-entity y de basketball; el segundo carece de evidencia confiable de entidad y produce una respuesta de incertidumbre. El rastreo propone un modelo contraintuitivo: no es que «IDK se active solo cuando algo es desconocido», sino que la tendencia IDK/refusal existe por defecto y las features de reconocimiento de entidades conocidas la inhiben.

\lecturefigure{slide-45-always-on-idk-circuit.jpg}{Las entidades conocidas suprimen mediante inhibition la ruta IDK predeterminada}{00:59:59}

\lecturefigure{slide-46-knowing-unknowns.jpg}{Comparación de las salidas para entidades conocidas y desconocidas}{01:00:17}

Lo representamos con un modelo de logits mínimo:
\[
z_{\text{answer}}=b_{\text{answer}}+s_{\text{knowledge}},\qquad
z_{\text{IDK}}=b_{\text{IDK}}-\gamma s_{\text{known}},
\]
donde $b$ es el sesgo predeterminado, $s_{\text{knowledge}}$ es la evidencia de una respuesta concreta, $s_{\text{known}}$ es la evidencia de que «la entidad es reconocible» y $\gamma>0$ indica la intensidad de la inhibición. La red real es más compleja; esta expresión solo describe la relación de competencia según la cual «reconocer algo como conocido puede reducir IDK».

\lecturefigure{slide-47-suppress-idk-hallucinates.jpg}{Al suprimir artificialmente el circuito IDK, la entidad desconocida recibe por fuerza conjeturas inestables}{01:02:22}

La intervención es decisiva: al reducir la feature IDK, el modelo da conjeturas distintas sobre Michael Batkin, como chess, tennis o hockey. Esto muestra que el circuito tiene un efecto causal sobre el rechazo a responder, y también que «hacer que el modelo responda siempre» convierte la incertidumbre en hechos de apariencia concreta.

\lecturefigure{slide-48-natural-hallucination-hypothesis.jpg}{Las alucinaciones naturales podrían deberse a que una señal known-entity errónea suprime IDK}{01:03:22}

La figura 48 extiende el mecanismo a nombres reales y títulos de artículos: algunas features de familiaridad superficial podrían emitir por error la señal de «conozco esta entidad», con lo que se levanta la inhibición de IDK; como la evidencia de una respuesta concreta es insuficiente, el modelo completa con un título relacionado pero incorrecto. La formulación aquí debe ser «some hallucinations might»; no se pueden atribuir todos los errores a un único circuito IDK.

Las tres figuras de mecanismo siguen un orden argumentativo estricto. La comparación entre Michael Jordan y Michael Batkin plantea primero la distinción «conocido/desconocido»; el experimento de supresión de IDK muestra que la dirección de rechazo tiene un efecto causal sobre la conducta; solo el caso de los títulos de artículos la propone como explicación candidata de los errores naturales. El tercer paso no manipula directamente todo el conocimiento del mundo real, por lo que su conclusión es más débil que la del segundo. Unos buenos apuntes deben separar por niveles lo «observado», lo «respaldado por intervención» y lo «extrapolado por conjetura»; de lo contrario, el lector tomará la hipótesis como si ya se hubiera localizado la causa raíz de todas las alucinaciones.

\begin{warningbox}{Las alucinaciones no son un mecanismo único ni siempre un solo token erróneo}
El recuerdo de hechos, la coherencia en textos largos, la mala interpretación de resultados de herramientas, los conflictos entre instrucciones y la continuación creativa pueden llamarse hallucination. Este caso explica la supresión del rechazo en un tipo de preguntas sobre entidades; no ofrece un detector universal ni define los límites de la alucinación en todos los escenarios.
\end{warningbox}

\subsection{El jailbreak no es una falla del interruptor de seguridad, sino un relevo entre rutas en competencia}

Preguntar directamente «cómo fabricar una bomba» activa de inmediato los circuitos harmful-request y refusal; el acrostic prompt pide primero extraer las iniciales, el modelo produce BOMB dentro de esa tarea local y después la ruta de continuation lo empuja a seguir explicando. El modelo puede haber reconocido ya el peligro y, aun así, no detenerse a tiempo porque la intensidad de las features y el enrutamiento difieren entre etapas.

\lecturefigure{slide-49-jailbreaks-intro.jpg}{Parallel Processing: Jailbreaks}{01:03:47}

\lecturefigure{slide-50-jailbreak-comparison.jpg}{Comparación de la conducta ante una solicitud peligrosa directa y ante un jailbreak por iniciales}{01:03:55}

\lecturefigure{slide-51-default-refusal-circuit.jpg}{Ruta predeterminada: el concepto bomb impulsa harmful request y refusal}{01:04:28}

\lecturefigure{slide-52-jailbreak-keeps-going.jpg}{Pregunta central: si el modelo ya escribió BOMB, ¿por qué sigue generando?}{01:04:30}

Al leer este grupo de figuras hay que comparar la secuencia temporal de los tokens: en la etapa del prompt dominan las features acrostic/first-letter; después de generar BOMB, el concepto peligroso y refusal empiezan a reforzarse, pero la continuation lingüística y la ruta de «responder a la tarea del usuario» ya han adquirido inercia. Si un sistema de seguridad clasifica solo una vez al inicio del prompt, puede pasar por alto estados de riesgo que surgen más tarde durante la generación.

Esto también explica por qué «el modelo sabe que es peligroso y aun así continúa» no es una contradicción lógica. Una red neuronal puede contener a la vez evidencia de reconocimiento del peligro y evidencia para seguir respondiendo; el token final lo decide el logit relativo, y reconocer el riesgo no garantiza que la ruta refusal domine en cada instante. Por lo tanto, la mejora en el nivel del mecanismo no consiste solo en reforzar un classifier, sino en lograr que las features de riesgo influyan de forma sostenida en la decodificación en la posición y el momento correctos, y en permitir reenrutar o detener la generación cuando cambia su estado.

\begin{importantbox}{Conexión con la ingeniería de seguridad: llevar el monitoreo al proceso de generación}
Las figuras de mecanismo sugieren tres direcciones prácticas: monitorear de forma continua los estados intermedios, en lugar de clasificar solo la entrada; permitir que el modelo reevalúe y se detenga a mitad de la generación; y tratar «no sé/rechazar» como una conducta en competencia que debe calibrarse, no como algo que simplemente se prohíbe o se impone.
\end{importantbox}

\teachervoice{00:58:54--01:04:45, el docente describe tanto la hallucination como el jailbreak como rutas en paralelo: el reconocimiento, la respuesta, IDK, refusal y continuation pueden coexistir. Nota de clase: las conclusiones sobre mecanismos son locales; no se puede extender un diagrama de circuito atractivo hasta convertirlo en una teoría de seguridad completa.}

\subsection{Resumen de la sección}

Los casos de alucinación y de jailbreak amplían el «paralelismo» de la capacidad aritmética a los conflictos de conducta. La salida del modelo no es la decisión de un único módulo, sino la combinación de varias rutas de evidencia y de inhibición. La intervención puede demostrar que ciertas rutas tienen un efecto causal, pero sigue siendo necesario distinguir entre el mecanismo de un caso, el riesgo a nivel de producto y las conclusiones generales de seguridad.
